%  -----------------------------------------------------------------------------
%  Author         : Bimalka Piyaruwan Thalagala
%  GitHub         : https://github.com/bimalka98
%  -----------------------------------------------------------------------------

\documentclass[a4paper,11pt]{article}%,twocolumn
%% packages

\usepackage{blindtext} % needed for creating dummy text passages
%\usepackage{ngerman} % needed for German default language
\usepackage{amsmath} % needed for command eqref
\usepackage{amssymb} % needed for math fonts
\usepackage[colorlinks=true,breaklinks]{hyperref} % needed for creating hyperlinks in the document, the option colorlinks=true gets rid of the awful boxes, breaklinks breaks lonkg links (list of figures), and ngerman sets everything for german as default hyperlinks language
\usepackage[hyphenbreaks]{breakurl} % ben�tigt f�r das Brechen von URLs in Literaturreferenzen, hyphenbreaks auch bei links, die �ber eine Seite gehen (mit hyphenation).
\usepackage{xcolor}
\definecolor{c1}{rgb}{0,0,1} % blue
\definecolor{c2}{rgb}{0,0.3,0.9} % light blue
\definecolor{c3}{rgb}{0.3,0,0.9} % red blue
\hypersetup{
    linkcolor={c1}, % internal links
    citecolor={c2}, % citations
    urlcolor={c3} % external links/urls
}
%\usepackage{cite} % needed for cite
\usepackage[square,authoryear]{natbib} % needed for cite and abbrvnat bibliography style
\usepackage[nottoc]{tocbibind} % needed for displaying bibliography and other in the table of contents
\usepackage{graphicx} % needed for \includegraphics 
\usepackage{longtable} % needed for long tables over pages
\usepackage{bigstrut} % needed for the command \bigstrut
\usepackage{enumerate} % needed for some options in enumerate
%\usepackage{todonotes} % needed for todos
\usepackage{makeidx} % needed for creating an index
\makeindex
\usepackage{gensymb}
\usepackage{url}
\usepackage{psfrag}
\usepackage{multirow}
\usepackage{subfigure}
%% page settings

\usepackage[top=15mm, bottom=15mm,left=20mm,right=20mm]{geometry} % needed for page border settings
\parindent=0mm % for space of first line of new text block
\sloppy % for writing with hyphenless justification (tries to)
\hyphenation{} % use hyphenation of tolerance parametershttp://www.jr-x.de/publikationen/latex/tipps/zeilenumbruch.html
\hyphenpenalty=10000
\exhyphenpenalty=10000
\usepackage{fancyhdr} % needed for head and foot options
%% my macros

%% Text fomats
\newcommand{\tbi}[1]{\textbf{\textit{#1}}}

%% Math fonts
\newcommand{\bbA}{\mathbb{A}}
\newcommand{\bbB}{\mathbb{B}}
\newcommand{\bbC}{\mathbb{C}}
\newcommand{\bbD}{\mathbb{D}}
\newcommand{\bbE}{\mathbb{E}}
\newcommand{\bbF}{\mathbb{F}}
\newcommand{\bbG}{\mathbb{G}}
\newcommand{\bbH}{\mathbb{H}}
\newcommand{\bbI}{\mathbb{I}}
\newcommand{\bbJ}{\mathbb{J}}
\newcommand{\bbK}{\mathbb{K}}
\newcommand{\bbL}{\mathbb{L}}
\newcommand{\bbM}{\mathbb{M}}
\newcommand{\bbN}{\mathbb{N}}
\newcommand{\bbO}{\mathbb{O}}
\newcommand{\bbP}{\mathbb{P}}
\newcommand{\bbQ}{\mathbb{Q}}
\newcommand{\bbR}{\mathbb{R}}
\newcommand{\bbS}{\mathbb{S}}
\newcommand{\bbT}{\mathbb{T}}
\newcommand{\bbU}{\mathbb{U}}
\newcommand{\bbV}{\mathbb{V}}
\newcommand{\bbW}{\mathbb{W}}
\newcommand{\bbX}{\mathbb{X}}
\newcommand{\bbY}{\mathbb{Y}}
\newcommand{\bbZ}{\mathbb{Z}}
\usepackage[ framed, numbered]{matlab-prettifier}%framed,%
\usepackage{listings}
\usepackage{pythonhighlight}
\usepackage{pdfpages}

\begin{document}

\begin{titlepage}
\center % Center everything on the page

%-------------------------------------------------------------------------------------
%	HEADING SECTIONS
%------------------------------------------------------------------------------------
\textbf{\large Department of Electronic and Telecommunication Engineering}\\[0.5cm]
\textbf{\Large University of Moratuwa, Sri Lanka}\\[1cm]
\textbf{\large EN 2053 - Communication Systems and Networks}\\[2cm]
\includegraphics[width=0.3\textwidth]{figures/uomlogo}\\[2cm]

	
%-------------------------------------------------------------------------------------
%	TITLE SECTION
%------------------------------------------------------------------------------------
\textbf{\Huge Modelling Propagation Losses and MANETs}\\[0.5cm]
\textbf{\Large Assignment 2}\\[5cm]


%----------------------------------------------------------------------------------------
%	MEMBERS SECTION
%----------------------------------------------------------------------------------------

\textbf{\large Submitted by}\\[0.5cm]
\begin{minipage}{0.2\textwidth}
	\begin{flushleft}
		{\large Thalagala B.P.}\\[4mm]
		{\large Sauranga H.W.C.}\\[4mm]
		
		
	\end{flushleft}
\end{minipage}
\hspace{5mm}
\begin{minipage}{0.2\textwidth}
	\begin{flushright}
		{\large 180631J }\\[4mm]
		{\large 180574K }\\[4mm]
		
	\end{flushright}
\end{minipage}\\[1.5cm]

%----------------------------------------------------------------------------------------
%	DATE SECTION
%----------------------------------------------------------------------------------------
\textbf{\large Submitted on}\\[0.5cm]
\textbf{\Large \today} % Date, change the \today to a set date if you want to be precise

%----------------------------------------------------------------------------------------

\vfill % Fill the rest of the page with whitespace

\end{titlepage}

\tableofcontents
\listoffigures
\listoftables
\vfill


\begin{center}
	\textit{\textbf{Note:}}
\textit{{\tt MATLAB R2018a} of the MathWorks Inc. is used for the implementation.}
\end{center}
\pagebreak

%%-----------------------------------------------------------------------
\section{Task 1}

Please note that for the MATLAB implementation bit rate(bits/second) of the generator was assumed to be 10. As we consider BPSK for the Task 1 and 2, the symbol rate(symbols/second) is remain the same as the bit rate of the generator.

\subsection{Generation of an Impulse Train Representing BPSK Symbols}
Binary data of the generator $D \in \{0, 1\}$ is mapped in to an impulse train according to the following function where $A(Amplitude)$  of the impulse was taken as 1 in the MATLAB implementation.
\[
amplitude ~of~ the~ k^{th}~ impulse = \begin{cases}
	+A & if~ D = 1\\
	-A & if~ D = 0
\end{cases}
\]

\subsection{Transmit Signal}
\subsection{Sinc function as the Impulse response}

Function {\tt sinc(t)} in MATLAB is defined as follows. 
\[
{\tt sinc(t)} = \begin{cases}
	\frac{sin(\pi.t)}{\pi.t} & if~t \neq 0\\
	1 & if~ t = 1
\end{cases}
\]

In order to generate  a sinc pulse that aligns with our time scale the function argument should be given as mentioned below. Where $T_b$ is the separation between
successive transmitted pulses.

\[
sinc~pulse = {\tt sinc(\frac{t}{T_b})}
\]


\subsection{Eye Diagram}

\section{Task 2}
Task 2 is a repetition of the Task 1, but in the presence of additive white Gaussian noise (AWGN). Variance of the noise $N_0/2$ was set such that the Power efficiency $\gamma = E_b/N_0 = 10 ~dB$. Where $E_b$ is the average bit energy and $N_0$ is the noise power spectral density.

\[
\begin{split}
	\gamma~in~dB &= 10.\log_{10}(E_b/N_0)\\
	\gamma/10 &= \log_{10}(E_b/N_0)\\
	10^{\gamma/10} &= E_b/N_0\\
	N_0 &= \frac{E_b}{10^{\gamma/10}}
\end{split}
\]

\[\therefore ~\sigma_{noise}^2 = N_0/2 = \frac{E_b}{2 \times 10^{\gamma/10}}  \]
Assuming the $P(D = 0) = P(D = 1) = 1/2$ as we consider sufficient large amount of bits in the initial bit stream,
\[
E_b = \frac{1}{2} \times \left[ (+A)^2 + (-A)^2 \right]
 = \frac{1}{2} \times \left[ (+1)^2 + (-1)^2 \right]
 = 1
\]
%---------------------------------------------------------------------------
\end{document}
